\section{Conclusion}
\label{sec:conclusion}

We address the trade-off between compression and cross-lingual parity in subword
tokenization by treating the per-language token costs as a vector of objectives. Its
power-mean scalarisation yields a greedy merge rule that weights each language's pair counts
by its current cost raised to the power $p-1$, which contains classical BPE at $p=1$ and
Parity-aware BPE as $p\to\infty$. Since pair scores never increase under this rule, the
resulting $p$-fair BPE trains exactly by lazy evaluation, about 20 times faster than full
rescoring. On Hindi, Marathi, Gujarati and English, raising $p$ traces a monotone and
repeatable front, on which a 40\% reduction in the Gini coefficient costs 0.6\% of
compression. The front coincides with that of the hybrid schedule, so the contribution is a
principled formulation of the existing trade-off rather than a tokenizer that dominates it.
Nevertheless, the formulation exposes a difference in how cost is distributed, since
% TODO (E3 in experiments.md): this "no language worse off" statement holds for the four-language
% setting only. In the 14-language pilot (reference_results/l14_preview) p-fair raises the cheapest
% languages (Bengali, Hindi) like every fairness rule does; what stays true is that p-fair never
% raises a language that costs MORE than the classical mean, whereas the hybrid does (Gujarati,
% Malayalam). Reword to the under-served form once E3 has error bars (paper_tables.py underserved).
$p$-fair BPE makes no Indian language more expensive than under classical BPE. Moreover, our
analysis shows that the cost unit and the training data balance decide the size, and even
the direction, of the disparity more than the merge rule does.

\noindent\textbf{Practical takeaways.} First, measure the cost on parallel content and not
per byte, since a per-byte cost inverts the parity rule across scripts. Second, report the
data balance alongside any parity result, since the same method removes 1.1\% or 7.0\% of
compression depending on how much English is present. Third, choose $p$ by the acceptable
loss in compression, since most of the disparity is removed well before full parity.

\noindent\textbf{Open directions.} The formulation extends naturally to further axes and
settings: a morphology penalty within the same rule (Section~\ref{sec:morph}), a convex
allocation of the vocabulary across scripts (Section~\ref{sec:alloc}), large vocabularies
where Parity-aware BPE loses equity \citep{lee2026}, and a $p$-fair variant of Unigram
pruning \citep{kudo2018}.

\section*{Limitations}

\noindent\textbf{Scale and languages.} We train on 140 million characters and evaluate on
1012 parallel sentences in four languages, three of which are closely related. The
sensitivity runs in Table~\ref{tab:sens} are single runs. Results may differ for larger
corpora, more languages or more distant scripts.

\noindent\textbf{Data balance as a design choice.} The size of the disparity depends mainly
on how much English is in the training text, and our main setting with four times as much
English is an assumption. Our conclusions should be read together with
Table~\ref{tab:sens}.

\noindent\textbf{Greedy, first-order steps.} Proposition~\ref{prop:rule} is a first-order
argument, and it assumes that pair frequencies in the development set match those in the
training text. Greedy merging carries no optimality guarantee for $p>1$, so the measured
curve is the front of the greedy algorithm and not the Pareto front of Eq.~\eqref{eq:obj}.

\noindent\textbf{Re-implemented baselines.} Parity-aware BPE, the hybrid schedule, the
moving window and Worst-$N$ are our own character-level implementations of the published
rules, whereas the originals are byte-level.

\noindent\textbf{Intrinsic evaluation.} We measure token cost and compression only. Lower
token cost does not by itself guarantee better downstream performance, and a model-level
evaluation is required to confirm that the gains matter in use.